% sections/10_appendix.tex

\section*{Appendix}
\label{sec:appendix}

\subsection*{A: Contract Schemas}
\label{app:contracts}

Full definitions of the 18 domain contracts and 12 request/response wrapper
dataclasses (30 frozen dataclasses total) that constitute the typed contract layer.
Each entry lists field names, types, default values, and validation
constraints enforced at instantiation.
\TODO{Insert auto-generated schema table from \texttt{chemometrics\_mcp.core.contracts}.}

\subsection*{B: Tool Signatures and JSON Examples}
\label{app:tool_signatures}

MCP-facing signatures (tool name, input schema, output schema) and
representative JSON request/response pairs for each of the 8 core analysis
tools and 3 method-memory tools.
\TODO{Insert tool signature table and JSON examples from server introspection.}

\subsection*{C: Per-Scenario Dataset Statistics}
\label{app:dataset_stats}

Sample counts, group counts, replicate counts per group, class distribution,
spectral axis range, and instrument metadata for each of the four evaluation
scenarios (S1: NIR wear-layer, S2: NIR species, S3: NIR vinyl/wood,
S4: FTIR purity).
\TODO{Insert dataset statistics table generated from the harness data loader.}

\subsection*{D: Ablation Configurations}
\label{app:ablation}

Full specification of each ablation configuration used in
Section~\ref{sec:ablation}: which components were disabled, what Python
flags or config overrides were applied, and which scenarios were run.
\TODO{Insert ablation configuration table from \texttt{evaluation/ablations/}.}

\subsection*{E: Fixture Construction}
\label{app:fixtures}

Description of how Layer-1 test fixtures were constructed for leakage
detection (positive and clean), hallucination probes (6 hand-authored),
HITL compliance checks, and fallback correctness.
Includes the criteria used to label each fixture as positive or clean,
and the validation procedure applied to confirm label correctness.
\TODO{Insert fixture inventory and construction rationale from
\texttt{evaluation/fixtures/README.md}.}

\subsection*{F: Reproducibility (Seeds, Versions, One-Command Run)}
\label{app:reproducibility}

Software versions (Python, scikit-learn, numpy, MCP SDK) pinned in
\texttt{requirements.lock}; random seed protocol (seed = 42 set globally
via \texttt{numpy.random.seed} and \texttt{random.seed} before each
scenario); and the per-scenario command to reproduce all Layer-1 results
(run once per scenario s1..s4):
\begin{verbatim}
python scripts/run_evaluation.py --scenario s1 --seed 42
\end{verbatim}
\TODO{Insert version table (Python, scikit-learn, numpy, MCP SDK).}

\subsection*{G: Expert Baseline Scripts}
\label{app:expert_baselines}

Scripts used to generate the expert baseline results reported in
Section~\ref{sec:evaluation}, including preprocessing choices, model
family, hyperparameter defaults, and CV strategy applied by the human
practitioner for each scenario.
\TODO{Insert baseline scripts from \texttt{baselines/} directory.}

\subsection*{H: Example Report --- Antigravity-Final Run}
\label{app:example_report}

The full \texttt{report.md} artifact generated by the
\texttt{antigravity-ftir-purity-run-final} case study described in
Section~\ref{sec:casestudy}, including the plan, preprocessing sweep
results, failed model records, validation warnings, and HITL decision
transcript.
\TODO{Insert content from
\texttt{runs/antigravity-ftir-purity-run-final/artifacts/report.md}.}
